\begin{abstract}
We prove the bounded-data Fourier restriction conjecture for the sphere in
three dimensions: the Fourier extension operator maps $L^\infty(S^2)$
boundedly into $L^p(\R^3)$ for every $p>3$. This is the full conjectured
open range, and the threshold $p=3$ is sharp.
\end{abstract}

\section{Introduction}
\label{sec:introduction}

Let \(S^2=\{\omega\in\R^3:|\omega|=1\}\), and let \(\sigma\) be surface
area measure, with \(\sigma(S^2)=4\pi\). For a bounded measurable
function \(g:S^2\to\C\), define
\[
  Eg(x)=\int_{S^2}g(\omega)e^{2\pi i x\cdot\omega}\,d\sigma(\omega).
\]
The bounded-data restriction problem asks for the exponents $p$ for which
this extension operator maps $L^\infty(S^2)$ boundedly into
$L^p(\R^3)$. We prove the full conjectured open range.

\begin{theorem}\label{thm:main}
For every real \(p>3\), there is a finite constant \(C_p\), depending
only on \(p\), such that
\[
  \|Eg\|_{L^p(\R^3)}
  \le C_p\|g\|_{L^\infty(S^2,\sigma)}
\]
for every bounded measurable complex-valued function \(g\) on \(S^2\).
\end{theorem}

The exponent range is sharp. For the constant function \(g=1\),
direct integration gives
\[
  E1(x)=\frac{2\sin(2\pi |x|)}{|x|}\qquad (x\ne0).
\]
The radial integral of \(|E1|^p\) diverges for \(0<p\le3\).
Thus Theorem~\ref{thm:main} proves exactly the conjectured open range;
it makes no endpoint assertion.

\subsection{The restriction problem and earlier methods}

The restriction problem originates in Stein's work, and the
Tomas--Stein theory gives the classical $L^2(S^2)\to L^4(\mathbb R^3)$
extension estimate~\cite{Stein1979,Tomas1975,Tomas1979}.
The oscillatory-integral work of Carleson--Sj\"olin and H\"ormander
\cite{CarlesonSjolin1972,Hormander1973}, Fefferman's ball-multiplier
obstruction~\cite{Fefferman1971}, and C\'ordoba's geometric multiplier
estimates~\cite{Cordoba1975} established important parts of the
analytic and geometric background. Bourgain's connection between
restriction and tube geometry~\cite{Bourgain1991Besicovitch} made
concentration of wave packets a central issue: curvature separates
packet directions, but many tubes can still pass through a small
region of space.

The bilinear framework of Tao--Vargas--Vega
\cite{TaoVargasVega1998} isolates interactions between separated
frequency pieces and reassembles them across separation scales by
rescaling and almost orthogonality. Wolff's cone theorem
\cite{Wolff2001} and Tao's paraboloid theorem~\cite{Tao2003}
developed this approach through wave packets and induction on scale.
In Tao's argument, a constraint from the Fourier analysis restricts
the directions of the packets that must be counted geometrically;
this supplies the additional structure needed to adapt Wolff's
counting argument to the paraboloid.

The multilinear theorem of Bennett--Carbery--Tao
\cite{BennettCarberyTao2006} controls interactions among surface
pieces whose normals uniformly span the ambient space. In three
dimensions it bounds the geometric mean of three extensions in
\(L^3\) on a ball, using the \(L^2\) norms of the three inputs and
an arbitrarily small positive power of the radius.
Bourgain--Guth~\cite{BourgainGuth2011} made this useful for the
linear problem by separating transverse interactions from
concentration in one small frequency region and from coplanar
interactions. Their argument treats the latter configurations by
rescaling and square-function estimates. This progression makes
the concentration of packet directions, as well as their
transverse intersections, a central part of the restriction problem.

Guth's polynomial-partitioning argument~\cite{Guth2016} divides the
extension integral between cells and a neighborhood of an algebraic
surface. Away from that neighborhood, each packet tube meets few of the
trimmed cells, allowing induction; the remaining concentration is
organized by tubes tangent to the surface. This gives the bounded-data
estimate for $p>13/4$ on a compact smooth positively curved surface $S$.
Shayya's weighted estimates~\cite{Shayya2017} give estimates with finite
$L^q(S)$ input norms in this output range. Kim~\cite{Kim2017} treats
these finite-input estimates through rescaling and interpolation. Wang's broom
method~\cite{Wang2022Brooms} then exploits a further constraint on the
truncated paraboloid: tubes that collect near one region spread toward
others, which limits repeated concentration near separated surface
neighborhoods. Combined with polynomial partitioning and a two-ends
argument, it yields the bounded-data estimate for $p>42/13$.

Wang and Wu~\cite{WangWu2022,WangWu2024} combine incidence estimates
with refined decoupling. The latter controls oscillatory interaction
using the number of packets meeting each small ball, while the incidence
argument retains how the occupied part of a tube is distributed along
its length. Their two-ends Furstenberg inequalities use this spacing
information to improve the multiplicity bound. In~\cite[Theorem~0.2]{WangWu2024}
they obtain the diagonal estimate
$L^p(S)\to L^p(\mathbb R^3)$ for $p>22/7$ on compact positively curved
$C^2$ surfaces, including boundary. The input norm matters in this
comparison: a diagonal estimate controls all $L^p(S)$ data, whereas
Theorem~\ref{thm:main} is stated for bounded sphere data. The latter
reaches the full open output range $p>3$; its sphere-specific
factorization consequence is discussed separately below.

\subsection{Elliptic capacity and the geometric obstacle}

The geometric and oscillatory estimates use the same eccentricity-biased
weight. For an ellipse \(E\subset\R^2\) with semiaxes \(L\ge w>0\), set
\[
 W(E)=L^{1+\kappa}w^{1-\kappa},\qquad 0<\kappa<1/10.
\]
A joint test for two vector coordinates uses translates of the same
ellipse, with the two centers chosen independently. It therefore
measures concentration of their joint law without assuming independence.
Equivalently, $W(E)=Lw(L/w)^\kappa$: the usual area weight is
multiplied by a power of the eccentricity. Differently oriented tests
intersect in smaller rectangles. In the hairbrush estimate of
Lemma~\ref{geo:hairbrush}, the extra exponent $\kappa$ turns the
resulting overlap count into a power saving. The propagation argument
uses that saving to control how concentration can change as a line moves.

Let \(M\) be a large dyadic scale, and consider a joint law of uniformly
bounded line parameters \((U,V)\) and a time \(T\in[0,1)\), with
\(X_t=V+tU\). Suppose its initial joint tests obey
\(\Pr(U\in u+E,V\in v+E)\le K W(E)\). The time hypothesis is
conditional on a base label that determines \(U,V\): at interval length
\(M^{-y}\), for \(0\le y\le1\), both the number of occupied time bins and the largest bin
probability have upper bounds with exponents \(sy\) and \(-sy\),
respectively, up to a small error. Here \(s\in(0,1]\).

At the terminal resolution let \(T_1\) be the left endpoint of the
time bin of length \(M^{-1}\) containing \(T\). A propagation exponent
\(q\) means that a dominated sublaw \(\nu\) satisfies
\[
 \nu(T_1=t,\ U\in u+E,\ X_t\in x+M^{-1}E)
       \le M^{-q+o(1)}K W(E)
\]
simultaneously for all terminal bins and independent centers \(u,x\).
The retained mass is \(M^{-o(1)}\), meaning that it exceeds every fixed
negative power for sufficiently large \(M\). Theorem~\ref{geo:main}
supplies every strict exponent below \(1+2s-10\kappa\) when the time
error is sufficiently small. That theorem specifies the input and
terminal ellipse scales and all conditional quantifiers.

\paragraph{Inputs and methodological precedents.}
The planar input belongs to the projection and Furstenberg-set
tradition of Marstrand and Kaufman~\cite{Marstrand1954,Kaufman1968},
Bourgain's discretized projection theorem~\cite{Bourgain2010}, and
Oberlin's work on exceptional projections~\cite{Oberlin2012}.
Packing-dimension advances of Orponen and Shmerkin
\cite{Orponen2020Packing,Shmerkin2022Packing}, their Hausdorff-dimension
improvement~\cite{OrponenShmerkin2023Hausdorff}, and their work relating
projections, Furstenberg sets, and the $ABC$ sum-product problem
\cite{OrponenShmerkin2026ABC} provide further context for Ren--Wang's
sharp finite incidence theorem~\cite{RenWang2023}. In the finite form
used here, each point of a planar family is incident to a comparably
sized family of thin tubes. Counts in smaller balls and angular
intervals control concentration of the points and tube directions.
Ren--Wang's theorem gives a lower bound for the total number of
distinct tubes. Drawing lines of constant projection through the
points converts this into a bound for projected supports. We prove
the weighted, conditional, and joint-coordinate consequences needed
for our laws.

The proof uses finite entropy in the sense of Shannon
\cite{Shannon1948} and the entropy rules organized by Tao
\cite{Tao2010Entropy}. The local-entropy approach of
Hochman--Shmerkin~\cite{HochmanShmerkin2012} and Hochman's inverse
principles~\cite{Hochman2014} are related methodological precedents;
our correlated conditional laws are handled by the finite arguments
below. The radial step is closely related to the reciprocal thin-tube
bootstrap of Orponen--Shmerkin--Wang
\cite{OrponenShmerkinWang2024Radial}. Hairbrush intersection counting
\cite{Wolff1995Kakeya,TaoVargasVega1998} and the density and factoring
arguments of Wang--Zahl and Guth--Wang--Zahl
\cite{WangZahl2025Convex,GuthWangZahl2026} supply further geometric
context. The independently translated elliptic capacity and its
propagation through arbitrary finite complex arrays and coordinate
masks are the specific estimates established here. The external
analytic input is Bourgain--Demeter's $\ell^2$ decoupling theorem
\cite{BourgainDemeter2015}. Section~\ref{sec:packets} proves the
canonical-packet refinement needed here, following the grouping and
rescaling method of Guth--Iosevich--Ou--Wang~\cite{GIOW2020} and
the localized extension formulation of Wang--Wu~\cite{WangWu2024}.
Besides these decoupling and finite-incidence inputs, the proof uses
classical almost-everywhere differentiation
\cite{Tao2011Measure,Heinonen2005Lipschitz}.

\subsection{The finite packet theorem}

The packet estimate underlying the sphere theorem applies to arbitrary
finite complex coefficient arrays. We summarize its definitions before
stating the uniform bound; Section~\ref{sec:packets} develops the frame
identities and norm properties in detail.

On a graph patch of the sphere, absorbing the surface-area density
into the data gives the phase \(x\cdot\xi+t\phi(\xi)\), with
\(\xi\in\R^2\). Rescaling a physical region of size \(N^2\) gives
the factor \(e^{2\pi iN^2(x\cdot\xi+t\phi(\xi))}\). A packet near
frequency \(\nu\), centered at \(z\) at time \(t_I\), is concentrated
near the line
\[
 x=z+(t-t_I)U_\nu,\qquad U_\nu=-\nabla\phi(\nu).
\]
Thus its frequency and position labels specify the velocity and
center tested by the elliptic capacity. The phase is continued
smoothly beyond the patch as specified in the theorem below.

For a fixed dyadic \(N\), a dyadic factor \(q\in[1,N]\) determines the
frequency width \(\theta=q/N\), time-bin length \(r=q^{-2}\), and
spatial mesh parameter \(w=(qN)^{-1}\). Choose a fixed smooth,
nonnegative, compactly supported window \(\chi\) whose integer
translates satisfy \(\sum_{n\in\mathbb Z^2}\chi(\xi-n)^2=1\).
Let \(L_{\mathrm f}\) be a fixed side length containing its support.
A bin \(I\) of length \(r\), with left endpoint \(t_I\), has packet
labels \((\nu,z)\in\theta\mathbb Z^2\times
L_{\mathrm f}^{-1}w\mathbb Z^2\). With
\(\chi_\theta^\nu(\xi)=\chi((\xi-\nu)/\theta)\), the exact maps for
the fixed phase \(\phi\) in the theorem below are
\begin{align*}
 (\mathsf P_{q,I}f)_{\nu,z}
 &=\int f(\xi)\chi_\theta^\nu(\xi)
       e^{2\pi iN^2(z\cdot\xi+t_I\phi(\xi))}\,d\xi,\\
 \mathsf S_{q,I}c(\xi)
 &=(L_{\mathrm f}\theta)^{-2}\sum_{\nu,z}
 c_{\nu,z}\chi_\theta^\nu(\xi)
       e^{-2\pi iN^2(z\cdot\xi+t_I\phi(\xi))}.
\end{align*}
Analysis thus takes unnormalized Fourier coefficients in each window;
synthesis has the corresponding inverse Fourier-series factor.

For an array \(c=(c_v)\) in one bin, indexed by
\(v=(\nu_v,z_v)\), put
\begin{align*}
 K(c)&=\theta^2\sup_{E,u,z}\frac{1}{W(E)}
       \sum_{\substack{U_{\nu_v}\in u+E\\z_v\in z+rE}}|c_v|^2,\\
 G(c)^3&=w^2\Bigl(\sum_v|c_v|^2\Bigr)K(c)^{1/2},\qquad
 A(c)=\inf_{|c|\le\sum_{j=1}^k|a_j|}\sum_{j=1}^kG(a_j).
\end{align*}
The supremum uses ellipses with both semiaxes at least \(\theta\),
and independent centers \(u,z\); the infimum uses finite arrays and
coordinatewise domination. The squared coefficients place mass
\(|c_v|^2\) at the velocity--position pairs \((U_{\nu_v},z_v)\).
Thus \(K\) measures joint concentration,
\(G\) combines that concentration with squared coefficient mass, and
the atomic norm \(A\) allows finite decompositions into arrays of
controlled \(G\)-cost. Subscripts below specify the scale and bin.

\begin{theorem}[Uniform propagation of finite packet arrays]\label{thm:intro-packets}
Fix \(0<\kappa<1/10\), a smooth elliptic continuation \(\phi\) of a
spherical graph patch allowed in Section~\ref{sec:packets}, a bounded
frequency-label region \(\Omega\subset\R^2\), and the fixed window system
used in \eqref{pkt:analysis}--\eqref{pkt:synthesis}.  The continuation is
defined on \(\R^2\), has uniformly definite bounded Hessian, and has
bounded derivatives of every order at least two.

For every finite \(B\ge1\) and every \(\varepsilon>0\), there is a finite
constant \(C\) with the following property.  Let \(m,N\) be dyadic with
\(1\le m\le N\le m^B\), and let \(d\) be any finite complex array on the
initial packet grid, at factor \(1\).  For each interval \(I\) in the
partition of \([0,1)\) into intervals of length \(m^{-2}\), let \(M_I\)
be any coordinate projection retaining finitely many entries of the
factor-\(m\) packet grid.  Assume that every initial index in
\(\operatorname{supp}d\) and every retained terminal index has frequency
label \(\nu\in\Omega\) and position label \(|z|\le m^B\).  Set
\[
 (T_md)_I=M_I\,\mathsf P_{m,I}\mathsf S_{1,[0,1)}d.
\]
Write \(A_1,G_1\) for the atomic norm and capacity gauge at the initial
scale
\[
 (\theta,r,w)=(N^{-1},1,N^{-1}),
\]
and \(A_{m,I}\) for the same atomic norm in the terminal bin \(I\), at
scale
\[
 (\theta,r,w)=(m/N,m^{-2},(mN)^{-1}).
\]
Then
\[
 m^{-2}\sum_{\substack{I\subset[0,1)\\ |I|=m^{-2}}}
       A_{m,I}\bigl((T_md)_I\bigr)^3
 \le C m^{10\kappa+\varepsilon}A_1(d)^3
 \le C m^{10\kappa+\varepsilon}G_1(d)^3,
\]
where the sum is over the specified partition.  The constant is uniform
in \(m,N,d\) and the coordinate projections.  It may depend on
\(B,\kappa,\varepsilon,\phi,\Omega\), and the fixed windows.
\end{theorem}

The maps retain the complete transition matrices between the selected
coordinate sets.  Each projection deletes coordinates; it does not
delete selected summands within a retained coefficient.  The initial
array need not be the analysis array of a function.  Frequency labels
may lie anywhere in the fixed region \(\Omega\) of the chosen
continuation, including its transition annulus.  The phase and its
derivative bounds are fixed before \(m,N\) vary.

Theorem~\ref{pkt:main} proves this estimate, with the atomic-input
formulation supplied by Lemma~\ref{pkt:atomic}.  The scale subscripts
above only specify which of the norms from Section~\ref{sec:packets}
is being used.

\subsection{Consequences and related estimates}

For the sphere, Bourgain's factorization theorem
\cite{Bourgain1991Besicovitch}, in the weak-type and interpolation
form explained in~\cite[Section~1.1, (1.1)--(1.2)]{Buschenhenke2024},
upgrades the full family in Theorem~\ref{thm:main} to
\[
 \|Eg\|_{L^q(\R^3)}\lesssim_q\|g\|_{L^q(S^2)}
 \qquad(3<q<\infty).
\]
This controls finite-$L^q$ sphere data as well as bounded data.
Interpolation also gives the open mixed-norm range in
Section~\ref{sec:consequences}.
The factorization is sphere-specific; the general-surface application
below uses the packet theorem itself.

\paragraph{A Kakeya maximal consequence.}
The sphere conclusion gives another route to the strong maximal estimate
proved independently in~\cite[Theorem~1.1]{OpenAIKakeya2026}.
For \(a\in\R^3\), \(\omega\in S^2\), and \(0<\delta<1\), put
\begin{align*}
 T_\delta(a,\omega)
 &=\{a+t\omega+u:|t|\le\tfrac12,\ u\perp\omega,\ |u|\le\delta\},\\
 K_\delta f(\omega)
 &=\sup_{a\in\R^3}\frac1{\pi\delta^2}
                  \int_{T_\delta(a,\omega)}|f(x)|\,dx.
\end{align*}
Thus the tubes have unit length and volume \(\pi\delta^2\).

\begin{corollary}\label{cor:sphere-kakeya-maximal}
For every \(\varepsilon>0\) there is a finite \(C_\varepsilon\) such that,
for every \(0<\delta<1\) and \(f\in L^3(\R^3)\),
\[
 \|K_\delta f\|_{L^3(S^2,\sigma)}
 \le C_\varepsilon\delta^{-\varepsilon}\|f\|_{L^3(\R^3)}.
\]
\end{corollary}

Sphere factorization, weighted tube duality, and interpolation give
this corollary; the complete transfer is proved in
Section~\ref{sec:consequences}. The supremum over tube locations is
part of the assertion, so it controls translated tube averages with
no restriction on where the tubes meet. The scale loss is an arbitrary
positive power. In particular, the standard implication from maximal
estimates gives Hausdorff dimension $3$ for every set containing a
unit line segment in every direction~\cite[Theorem~2.12]{MattilaFourierNotes}.
That dimension conclusion was already proved by Wang and
Zahl~\cite[Theorem~1.1]{WangZahl2025Convex}; their result also gives
Minkowski dimension $3$.

\paragraph{The full packet theorem on other surfaces.}
The full array formulation also has a separate application.
The companion article~\cite[Theorem~1.1]{OpenAIDiagonal2026} combines it with a
translated-tube Kakeya maximal estimate to prove
$L^p(\Sigma)\to L^p(\mathbb R^3)$ extension for every $3<p<\infty$
on each compact smooth positively curved surface, possibly with smooth boundary.
For spheres and quadratic graphs, bounded-data-to-diagonal
factorization already has an established history
\cite{Bourgain1991Besicovitch,Buschenhenke2024}; the companion treats
general elliptic charts using the full packet theorem and an explicit
phase continuation. Recent smooth Alpert testing formulations
\cite{Sawyer2026TestingComparison,RiosSawyer2026ConvolutionTesting}
study extension through specified testing conditions. The conclusions
here concern deterministic complex arrays and all bounded measurable
sphere data, with the fixed phase, window, and complexity parameters
specified in the theorem.

\paragraph{An independent Bochner--Riesz route.}
The bounded-data spherical estimate in Theorem~\ref{thm:main} also
follows independently from the Bochner--Riesz multiplier theorem
in~\cite[Corollary~12.4]{OpenAIBochnerRiesz2026}, through Tao's
implication~\cite{Tao1999Epsilon}. This implication does not recover
the capacity and packet estimates proved here.

\subsection{The argument}

The eccentricity bias in \(W\) limits the mass that can concentrate
simultaneously in differently oriented tests. The proof develops this
fact for conditional laws, applies the resulting propagation theorem
to packet arrays, and finally removes the spatial power loss.

\paragraph{Geometric propagation.}
Suppose the retained-law bound above fails. Extremizing over relative
spatial scales selects tests whose mass is large compared with their
capacity weight at successive time resolutions. Writing their radii as
powers of \(M\) gives two logarithmic width-depth functions of time
depth. The hairbrush bound
controls differences between their axes, assigning each sample of
the joint law a reference axis on a short interval of scales.
Near a differentiability point of the two width-depth functions, their
rates describe whether the aspect ratio is stationary, increasing,
or decreasing. This is the calibrated width path constructed in
Section~\ref{sec:geo}.

We record conditional Shannon entropies of velocity, position, time,
and the reference axis, rounding these variables at compatible
resolutions. Time and axis may be correlated. The laws are first
regularized on finite grids, so the limiting entropy increments
retain the support counts and probability bounds required by the
projection theorems. Those theorems force lower bounds on entropy
growth. In each of the three aspect cases, the bounds contradict the
concentration required by failure of geometric propagation.

\paragraph{Propagation of wave-packet bounds.}
The exact frame maps compose between packet scales, and the atomic
norm lets us combine estimates despite intermediate coordinate
restrictions. Insert finitely many intermediate scales and form a
layered graph whose vertices are the retained packet indices. An edge
records that a child frequency and center lie in a parent's
localization neighborhood. A path is a sequence of such indices from
an initial packet to a terminal one; the coefficients still evolve
by the full masked matrices.

Suppose amplification has a growth exponent greater than
\(10\kappa\). Among families approaching the largest exponent, first
take the infimum \(s\) of the exponents \(\sigma\) for which each
initial index reaches at most \(m^{2\sigma+o(1)}\) terminal time bins.
This infimum is positive; it need not be attained.
Then minimize the number of paths per original initial index, again
at the level of scale exponents. The second extremum forces later
coordinate selections to retain a subpower fraction of those paths
whenever they preserve nearly maximal amplification. Uniform counting
on selected paths then gives both conditional prefix counts and
maximal prefix probabilities with exponent \(s\) in the geometric base
\(M=m^2\). The initial packet index is the base label and determines
the line parameters \((U_\nu,z)\). This is the law to which geometric
propagation applies.

The geometric bound controls output capacity. Refined decoupling
controls a fourth moment, while the frame identity gives a second
moment bound. Their combination cancels the possible time-support
deficit and leaves a growth exponent at most \(10\kappa\).
Figure~\ref{fig:packet-proof-map} shows how the two extrema lead to
the conditional time law and the capacity--mass cancellation.

\paragraph{Removing the spatial loss.}
The atomic norm controls the normalized sum of cubed packet
coefficients. For actual graph-patch data, terminal coefficients are
samples of the localized extension, up to rapidly decaying truncation
errors. Averaging the sampling lattice therefore gives a cubic
integral estimate. Taking \(\kappa\) small makes its spatial power
loss arbitrarily small. On widely separated balls, a joint frame
estimate bounds the total initial coefficient mass without a factor
for the number of balls.
A covering of the large-value set, followed by Fourier reproduction
and layer integration, yields the global estimate for every \(p>3\).
This last step works directly for bounded measurable complex data.

\paragraph{Structural statements within the proof.}
The geometric and entropy arguments also yield structural statements.
The matrix form of the ellipse weight
obeys the composition inequality \eqref{geo:submultiplicativity}.
Corollary~\ref{geo:axis-prediction} determines a sample's axis, up to
the corresponding aspect precision, from its exact line parameters
and time prefix. Lemmas~\ref{ent:peeling} and~\ref{ent:signed-envelope} describe
additive parts of the entropy's local limits and lower bounds when
some resolutions increase along a path and others decrease. They
also apply to the specified paths on boundary faces of the admissible
resolution domain.

\subsection{Organization and conventions}

Sections~\ref{sec:geo}--\ref{sec:aspect} prove the geometric
propagation theorem. Section~\ref{sec:geo} constructs the calibrated
path; Section~\ref{sec:entropy} develops its entropy profiles and
tangent rules; Section~\ref{sec:projection} proves the projection
consequences of the planar Furstenberg theorem; and
Section~\ref{sec:aspect} treats the three possible directions of
aspect change.
Sections~\ref{sec:packets} and~\ref{sec:extremal} prove the
wave-packet estimate, and Section~\ref{sec:global} concludes the proof
of Theorem~\ref{thm:main}. Section~\ref{sec:consequences} gives the
sphere-to-Kakeya transfer and the mixed-norm sphere estimates.

Functions may be complex-valued, and probability laws may have
arbitrary correlations, unless a more restrictive hypothesis is
explicitly stated. Ellipse containments suppress independent translations
when two spatial coordinates are tested. Fixed multiplicative
constants are harmless in scale exponents. The geometric and
packet sections specify their limiting conventions separately;
in each case finite grids and fixed complexity are chosen before
the large-scale limit. A \emph{subpower} loss denotes a factor
\(M^{o(1)}\), or \(m^{o(1)}\), in the base currently under discussion.
